\chapter{Conclusion}
\label{chap:conclusion}

In this work, a Data-Driven Min-Max Model Predictive Control (DD-MPC) strategy was designed, implemented, and experimentally validated for the stabilization of a self-balancing mini-bicycle.
Unlike traditional model-based methods, the data-driven approach leveraged collected input-state measurements directly, eliminating the need for an explicit mathematical model.

Through simulations and real-world experiments, the proposed Data-Driven MPC controller demonstrated improved performance compared to a baseline Linear Quadratic Regulator (LQR) controller.
Specifically, the data-driven controller achieved faster stabilization, lower overshoot, and better robustness against model uncertainties and external disturbances.

Despite hardware limitations, such as computational restrictions on the STM32 microcontroller and sensor noise, the controller effectively stabilized the mini-bicycle within defined safe operating limits.
The offline optimization approach combined with precomputed state-feedback gains enabled real-time embedded operation, confirming the feasibility of deploying advanced data-driven control strategies on resource-constrained platforms.

\section*{Future Work}

Future research directions include:
\begin{itemize}
	\item Development of an online adaptive Data-Driven MPC to update the controller in real-time based on newly collected data.
	\item Exploration of rear-wheel steering strategies for stationary balancing without forward motion.
	\item Integration of higher-performance embedded hardware to allow fully online SDP optimization onboard.
	\item Testing the controller under varying payload conditions and dynamic environments to further evaluate robustness.
	\item data driven Sys Identification
\end{itemize}

Future work could focus on several promising directions to extend this research. First, the development of an online adaptive data-driven MPC would enable the controller to update in real time based on newly collected data, enhancing adaptability to changing system dynamics. Additionally, exploring rear-wheel steering strategies could allow for stationary balancing without requiring forward motion, broadening the operational capabilities of the bicycle. Another key avenue is the integration of higher-performance embedded hardware, which would make it feasible to perform fully online SDP optimization directly onboard the bicycle. Furthermore, extensive testing under varying payload conditions and in dynamic environments would be essential to rigorously evaluate and improve the robustness of the controller. Finally, applying data-driven system identification techniques could complement the control design by offering insights into the underlying dynamics and refining controller performance.

